% ============================================================
% AEGIS unified physics report — "all the physics, in all cases"
% A single formal derivation of every physical model in the AEGIS
% absorbed-power-density engine: the incoherent surface law, the
% Fock diffraction gate (local angle metric and distal C-metric),
% impedance creeping waves, the coherent exposure operator, near
% and far field, inter-body reflection, and the psSAR10g bridge.
%
% Build:  cd theory && latexmk -pdf unified_physics.tex
%   (or:  cd theory && pdflatex unified_physics && pdflatex unified_physics)
% Reuses the monograph preamble (unified/_preamble_base.tex) so notation
% matches theory/monograph_v2.tex. Section bodies live in unified/sec_*.tex.
% ============================================================
\documentclass[11pt,a4paper]{article}

% ---------------------------------------------------------------
\usepackage[margin=2.4cm]{geometry}

% ---------------------------------------------------------------
% ENCODING AND FONTS
% ---------------------------------------------------------------
\usepackage[utf8]{inputenc}
\usepackage[T1]{fontenc}
\usepackage{lmodern}           % CM-compatible outlined fonts
% Alternative: Times-like math (closer to AMS / IEEE style)
%   \usepackage{newtxtext,newtxmath}  % requires newtx package
\usepackage{microtype}         % character-level justification; always last font pkg
\linespread{1.04}              % slight openness; 1.0 is pure single-spacing

% ---------------------------------------------------------------
% MATHEMATICS — core
% ---------------------------------------------------------------
\usepackage{amsmath,amssymb,amsthm,mathtools}
\usepackage{bm}                % bold math (vectors, matrices)
\usepackage{physics}           % \abs, \norm, \dd, \grad, …
\usepackage{mathrsfs}          % \mathscr{} script letters (e.g. for operators)
\usepackage{bbm}               % \mathbbm{1} blackboard-bold indicator
\usepackage{cancel}            % \cancel, \bcancel for struck-through terms
\usepackage{xfrac}             % \sfrac{}{} for inline slanted fractions

% Allow long equation arrays to break across pages
\allowdisplaybreaks[2]

% ---------------------------------------------------------------
% MATHEMATICS — highlighted / boxed equations
% ---------------------------------------------------------------
\usepackage{empheq}            % \begin{empheq}[box=\fbox]{align}
% Define a shaded box option for main results:
\usepackage[most]{tcolorbox}
\tcbset{highlight math style={%
    enhanced, colframe=black!70, colback=gray!6,
    arc=2pt, boxrule=0.5pt, left=3pt, right=3pt}}

% ---------------------------------------------------------------
% SI UNITS
% ---------------------------------------------------------------
\usepackage{siunitx}           % \SI{28}{\giga\hertz}, \si{\watt\per\kilo\gram}
\sisetup{
  detect-all,                  % inherit surrounding font
  range-phrase = {--},
  range-units  = single,
  per-mode     = symbol,
  inter-unit-product = {\,}
}

% ---------------------------------------------------------------
% FIGURES AND TABLES
% ---------------------------------------------------------------
\usepackage{graphicx}
\graphicspath{{figures/}}
\usepackage{booktabs}          % \toprule, \midrule, \bottomrule
\usepackage{array}             % extended column specifiers
\usepackage{tabularx}          % full-width tables
\usepackage{multirow}
\usepackage{makecell}          % multi-line table cells
\usepackage{caption}
\captionsetup{font=small, labelfont=bf, labelsep=period,
              skip=4pt}        % compact captions, bold label, period separator
\usepackage{subcaption}
\usepackage{threeparttable}    % table notes for lit-review tables
\usepackage{float}
\usepackage[section]{placeins} % keep floats within their section

% Relax float placement rules
\renewcommand{\topfraction}{0.9}
\renewcommand{\bottomfraction}{0.9}
\renewcommand{\textfraction}{0.1}
\renewcommand{\floatpagefraction}{0.7}
\setcounter{topnumber}{4}
\setcounter{bottomnumber}{4}
\setcounter{totalnumber}{8}

% ---------------------------------------------------------------
% SECTION HEADINGS (titlesec)
% ---------------------------------------------------------------
\usepackage{titlesec}
\titleformat{\section}{\large\bfseries}{\thesection}{1em}{}
\titleformat{\subsection}{\normalsize\bfseries}{\thesubsection}{1em}{}
\titleformat{\subsubsection}{\normalsize\bfseries}{\thesubsubsection}{1em}{}
\titlespacing*{\section}{0pt}{1.6ex plus .4ex minus .2ex}{0.8ex plus .2ex}
\titlespacing*{\subsection}{0pt}{1.2ex plus .3ex minus .2ex}{0.5ex plus .1ex}

% ---------------------------------------------------------------
% MISCELLANEOUS
% ---------------------------------------------------------------
\usepackage{enumitem}          % fine-grained list control
\usepackage{xcolor}
\usepackage{xspace}            % smart spacing after macro names
\usepackage{todonotes}[disable]% \todo{} comments — set [disable] for final version
\usepackage{fancyhdr}          % running headers

% Running header: empty header, page number centred in footer
\pagestyle{fancy}
\fancyhf{}
\cfoot{\thepage}
\renewcommand{\headrulewidth}{0pt}   % no rule when header is empty

% ---------------------------------------------------------------
% BIBLIOGRAPHY
% ---------------------------------------------------------------
\usepackage[numbers,sort&compress]{natbib}

% ---------------------------------------------------------------
% HYPERLINKS — load near last, before cleveref
% ---------------------------------------------------------------
\usepackage[
  colorlinks = true,
  linkcolor  = {blue!55!black},
  citecolor  = {green!45!black},
  urlcolor   = {blue!65!black},
  pdfstartview = FitH,
  bookmarksnumbered = true,
  bookmarksopen     = true
]{hyperref}
\usepackage{bookmark}          % fixes bookmark levels in PDF outline
\usepackage{doi}               % live DOI hyperlinks (must load AFTER hyperref)

% cleveref must come AFTER hyperref
\usepackage[nameinlink,capitalize]{cleveref}

% ---------------------------------------------------------------
% THEOREM-LIKE ENVIRONMENTS
% (all share the same counter as theorem, numbered per section)
% ---------------------------------------------------------------
\usepackage{aliascnt}

\theoremstyle{definition}      % upright body text for all env.

\newtheorem{theorem}{Theorem}[section]

\newaliascnt{proposition}{theorem}
\newtheorem{proposition}[proposition]{Proposition}
\aliascntresetthe{proposition}

\newaliascnt{lemma}{theorem}
\newtheorem{lemma}[lemma]{Lemma}
\aliascntresetthe{lemma}

\newaliascnt{corollary}{theorem}
\newtheorem{corollary}[corollary]{Corollary}
\aliascntresetthe{corollary}

\newaliascnt{definition}{theorem}
\newtheorem{definition}[definition]{Definition}
\aliascntresetthe{definition}

\newaliascnt{remark}{theorem}
\newtheorem{remark}[remark]{Remark}
\aliascntresetthe{remark}

% cleveref names for theorem-like environments
\crefname{proposition}{Proposition}{Propositions}
\Crefname{proposition}{Proposition}{Propositions}
\crefname{lemma}{Lemma}{Lemmas}
\Crefname{lemma}{Lemma}{Lemmas}
\crefname{corollary}{Corollary}{Corollaries}
\Crefname{corollary}{Corollary}{Corollaries}
\crefname{definition}{Definition}{Definitions}
\Crefname{definition}{Definition}{Definitions}
\crefname{remark}{Remark}{Remarks}
\Crefname{remark}{Remark}{Remarks}

% ---------------------------------------------------------------
% CUSTOM MATH COMMANDS
% ---------------------------------------------------------------
% --- vectors and fields ---
\newcommand{\khat}{\hat{\bm{k}}}
\newcommand{\nhat}{\hat{\bm{n}}}
\newcommand{\rr}{\mathbf{r}}
\newcommand{\EE}{\mathbf{E}}
\newcommand{\HH}{\mathbf{H}}
% --- dosimetric quantities ---
\newcommand{\Sinc}{S_{\mathrm{inc}}}
\newcommand{\Sab}{S_{\mathrm{ab}}}
\newcommand{\Teff}{T_{\mathrm{eff}}}
\newcommand{\Tavg}{T_{\mathrm{avg}}}
\newcommand{\Aab}{A_{\mathrm{ab}}}
\newcommand{\Aperp}{A_{\perp}}
\newcommand{\SARwb}{\mathrm{SAR}_{\mathrm{wb}}}
% --- material parameters ---
\newcommand{\ntilde}{\tilde{n}}
\newcommand{\epstilde}{\tilde{\varepsilon}_r}
% --- operators ---
\DeclareMathOperator{\ReLU}{ReLU}
% (\erf already provided by the physics package)
\DeclareMathOperator{\RE}{Re}
\DeclareMathOperator{\IM}{Im}
% --- other shorthands ---
\newcommand{\Tbar}{\bar{T}}
\newcommand{\diff}{\mathrm{d}}        % upright differential
\newcommand{\ident}{\mathbbm{1}}       % indicator / identity  (requires bbm)
\newcommand{\Reals}{\mathbb{R}}
\newcommand{\Complex}{\mathbb{C}}
\newcommand{\half}{\tfrac{1}{2}}

% ---------------------------------------------------------------
% PART HEADINGS (article class has no \part)
% ---------------------------------------------------------------

% ============================================================
% Additional macros for the unified physics report
% (diffraction, curvature, coherent). Loaded after the monograph base preamble.
% ============================================================
% --- Fock diffraction ---
\newcommand{\mfock}{m}                       % Fock large parameter (kR/2)^{1/3}
\newcommand{\Rfock}{R}                        % in-incidence-plane radius of curvature
\newcommand{\qF}{q_{F}}                       % Fock surface-impedance parameter
\newcommand{\Philit}{\Phi_{\mathrm{lit}}}     % lit transition
\newcommand{\Psishad}{\Psi_{\mathrm{sh}}}     % shadow creeping series
\newcommand{\Ai}{\operatorname{Ai}}           % Airy function
\newcommand{\Aip}{\operatorname{Ai}'}         % Airy derivative
\newcommand{\Rocc}{R_{\mathrm{occ}}}          % binding-occluder radius (distal)
\newcommand{\clr}{c}                          % signed angular clearance (C-metric)
% --- coherent ---
\newcommand{\Gtil}{\widetilde{\mathbf{G}}}    % body-surface channel
\newcommand{\Qop}{\mathbf{Q}}                 % exposure operator
\newcommand{\xx}{\mathbf{x}}                  % precoder vector
\newcommand{\psivec}{\bm{\psi}}               % polarisation-amplitude vector
\newcommand{\herm}{^{\mathsf{H}}}             % Hermitian transpose


\graphicspath{{figures/}{unified/figures/}}

\title{\textbf{AEGIS: A Unified Derivation of Geometric and\\[0.3em]
Diffractive Absorbed-Power-Density Dosimetry}}
\author{%
  Robin Wydaeghe \\
  Ghent University \& imec, Ghent, Belgium \\
  \normalfont\texttt{robin.wydaeghe@ugent.be}
}
\date{\today}

\makeatletter
\renewcommand{\maketitle}{%
  \thispagestyle{plain}%
  \begin{center}
    \rule{\linewidth}{1pt}\\[0.4em]
    {\LARGE\bfseries \@title}\\[0.5em]
    {\large \@author}\\[0.3em]
    {\small \@date}
    \\[0.4em]\rule{\linewidth}{1pt}
  \end{center}
  \vspace{0.5em}
}
\makeatother

\begin{document}
\maketitle

\begin{abstract}
This report derives, from first principles and in one consistent notation, the
complete physical model used by the AEGIS absorbed-power-density engine. We start
from the local surface law $\Sab(\rr)=\Sinc\,T_0\,\ReLU[\nhat\cdot(-\khat)]\,O$
and build outward: tissue electromagnetics and Fresnel transmission; the
orthogonal fidelity grid; polarisation-aware absorption; the geometric quantities
(projected area, occlusion, principal curvature); the Fock diffraction gate, in
both its \emph{local} terminator-angle form and its \emph{distal} angular-clearance
(``C-metric'') form, including the soft/hard impedance creeping-wave eigenvalues;
inter-body reflection; the near- and far-field limits; the coherent exposure
operator and energy-constrained beamforming; and the bridge to peak spatial SAR.
Every model is stated formally, with its domain of validity and its validation
against exact references (Mie spheres, Bessel--Hankel cylinders). The convention
throughout is $\ntilde = n - \mathrm{i}\kappa$ ($e^{+\mathrm{i}\omega t}$).
\end{abstract}

\tableofcontents
\bigskip

% ---- sections (authored in unified/sec_*.tex) ----
% sec_01_surface_law
\section{Setup and the master surface law}\label{sec:surface-law}

AEGIS reduces dosimetry to a single scalar field on the body surface: the
absorbed power density $\Sab(\rr)$, in watts per square metre, deposited at a
skin point $\rr$. Every fidelity level in this report computes $\Sab$, and
every refinement (Fresnel transmission, polarisation, curvature, diffraction,
inter-body coupling, coherent beamforming) is a correction to one factor of a
single master law. This section states that law, defines each symbol, and
fixes the sign and time conventions that the diffraction and coherent sections
depend on.

\subsection{Geometric setup}\label{sec:setup-geometry}

Consider a monochromatic plane wave of angular frequency $\omega = 2\pi f$
propagating in the unit direction $\khat$, with time-averaged Poynting vector
$\mathbf{S}_{\mathrm{inc}} = \Sinc\,\khat$. The scalar $\Sinc$ is the incident
power spectral density (the magnitude of the time-averaged Poynting flux,
\si{\watt\per\square\metre}). The body surface $\Sigma$ is described by an
outward unit normal field $\nhat(\rr)$ and is assumed locally flat on the scale
of the wavelength, so that a Fresnel half-space model applies pointwise. This
local-flatness assumption is the one relaxed by the curvature and diffraction
refinements of \cref{sec:grid} and the later geometry sections.

\begin{definition}[Incidence cosine]\label{def:mu}
At a surface point $\rr$ with outward normal $\nhat(\rr)$, the incidence cosine
is
\begin{equation}\label{eq:mu-def}
  \mu(\rr) \;\equiv\; \nhat(\rr)\cdot(-\khat) \;=\; \cos\theta_i(\rr),
\end{equation}
where $\theta_i$ is the local angle of incidence. A front-facing point has
$\mu > 0$, a point in geometric shadow has $\mu \le 0$.
\end{definition}

\subsection{Time-averaged power balance}\label{sec:power-balance}

The inward power flux through a surface element $\diff A$ at $\rr$ is
$\diff P_{\mathrm{in}} = \Sinc\,\mu\,\diff A$: the beam carries $\Sinc$ per unit
area normal to $\khat$, and a tilted element of its own area $\diff A$
intercepts only the projected fraction $\mu = \cos\theta_i$. The specularly
reflected wave leaves at the same angle with power density
$\abs{r(\theta_i)}^2\,\Sinc$, where $r(\theta_i)$ is the Fresnel amplitude
reflection coefficient for the relevant polarisation, so the outward reflected
flux is $\diff P_{\mathrm{ref}} = \abs{r}^2\,\Sinc\,\mu\,\diff A$.

At millimetre-wave frequencies the skin depth is $\delta \sim \SIrange{0.3}{0.5}{\milli\metre}$,
far below any body curvature radius, so all transmitted power is absorbed in
the surface layer. Energy conservation at the interface gives
$\diff P_{\mathrm{abs}} = \diff P_{\mathrm{in}} - \diff P_{\mathrm{ref}}
= \Sinc\,\mu\,(1-\abs{r}^2)\,\diff A$.

\begin{definition}[Power transmission coefficient]\label{def:T}
For a given polarisation the power-absorption (transmission) coefficient is
$T(\theta_i) \equiv 1 - \abs{r(\theta_i)}^2$. Its normal-incidence value is
written $T_0 \equiv T(0)$.
\end{definition}

\subsection{The master surface law}\label{sec:master-law}

Dividing the absorbed flux by $\diff A$ and gating out back-facing points
yields the law that organises the entire report.

\begin{proposition}[Master surface law]\label{prop:master-law}
The absorbed power density at a surface point $\rr$ is
\begin{equation}\label{eq:master-law}
  \boxed{\;\Sab(\rr) \;=\; \Sinc\;T_0\;\ReLU\!\bigl[\nhat(\rr)\cdot(-\khat)\bigr]\;
  O(\rr,\khat)\;}
\end{equation}
where $\ReLU(\mu)=\max(0,\mu)$ and $O(\rr,\khat)\in\{0,1\}$ in the
geometric-optics limit is the visibility factor, smoothed by the diffraction gate
in later sections. Each factor is independent and is refined separately in later
sections.
\end{proposition}

The five factors carry the whole physics:
\begin{itemize}[leftmargin=1.4em,itemsep=1pt]
  \item $\Sinc$: incident power spectral density at $\rr$
    (\si{\watt\per\square\metre}), supplied by the ray tracer or environment
    model.
  \item $T_0$: normal-incidence power transmission into tissue
    (\cref{def:T}), the material factor. \Cref{sec:tissue} derives it from the
    tissue refractive index and \cref{sec:fresnel} generalises it to the
    angle- and polarisation-dependent $\Teff(\rr)$.
  \item $\ReLU[\mu]$: the cosine projection combined with a back-face cull. It
    is zero in geometric shadow and equal to $\cos\theta_i$ on the lit side.
    \Cref{sec:grid} and the diffraction sections soften this hard gate into a
    physically derived smooth transition near the shadow boundary.
  \item $O(\rr,\khat)\in\{0,1\}$ in the geometric-optics limit: the visibility or occlusion factor, unity for a
    directly illuminated point and reduced by self-shadowing, inter-body
    occlusion, and ambient occlusion. It is developed in the geometry and
    inter-body sections.
\end{itemize}

For unpolarised or circularly polarised illumination the material factor is the
scalar $T_0$ shown in \cref{eq:master-law}. The exact, polarisation-resolved
form replaces $T_0$ with the effective transmission $\Teff(\rr)$ of
\cref{sec:pol}, giving
\begin{equation}\label{eq:master-law-exact}
  \Sab(\rr) \;=\; \Sinc\;\Teff(\rr)\;\ReLU\!\bigl[\mu(\rr)\bigr]\;O(\rr,\khat),
\end{equation}
which is exact wherever the surface is locally flat. The coherent levels of
\cref{sec:grid} replace the incoherent power product entirely with a
quadratic form in the complex precoder, recovering \cref{eq:master-law-exact}
in the incoherent limit.

\subsection{Sign and time conventions}\label{sec:conventions}

The following convention is load-bearing. The Fresnel branch cuts, the
penetration depth, and the Fock attenuation in the diffraction sections all
depend on it, so we state it once and use it everywhere.

\begin{remark}[Refractive-index and time convention]\label{rem:convention}
Fields vary in time as $e^{+\mathrm{i}\omega t}$. With this convention a lossy
medium has complex relative permittivity
$\epstilde = \varepsilon_r - \mathrm{i}\,\sigma/(\omega\varepsilon_0)$ and
complex refractive index
\begin{equation}\label{eq:n-convention}
  \ntilde \;=\; \sqrt{\epstilde} \;=\; n - \mathrm{i}\kappa,
  \qquad n>0,\;\kappa>0,
\end{equation}
where the branch is fixed by $\RE(\ntilde)>0$ and $\IM(\ntilde)<0$ so that a
wave decays as it penetrates the tissue. Throughout, $\mu = \nhat\cdot(-\khat)$
denotes the incidence cosine of \cref{def:mu}. Adopting the opposite
($e^{-\mathrm{i}\omega t}$) convention flips the sign of $\kappa$ and of every
imaginary part below.
\end{remark}

% sec_02_tissue
\section{Tissue electromagnetics}\label{sec:tissue}

\emph{[section pending]}

% sec_03_fresnel
\section{Fresnel transmission}\label{sec:fresnel}

\emph{[section pending]}

% sec_04_fidelity_grid
\section{The fidelity grid}\label{sec:grid}

AEGIS exposes nine fidelity levels, numbered 0 to 8, that trade accuracy for
cost. Rather than nine unrelated models, they are organised as an orthogonal
composition: a small set of physical corrections, each touching a different
structural element of the master law, switched on independently. This section
states the composition, lists the levels, and identifies the new orthogonal
axes (diffraction model and inter-body coupling) developed later in the report.

\subsection{Two regimes}\label{sec:two-regimes}

The levels split into two regimes by what they act on.

The incoherent levels (0 to 6) act on per-path power. Each propagation path
contributes a scalar power density $S_i$, and absorbed power densities add as
powers,
\begin{equation}\label{eq:incoherent-sum}
  \Sab(\rr) \;=\; \sum_i S_i\;\Teff^{(i)}(\rr)\;g\!\bigl(\mu_i(\rr),\sigma\bigr)\;
  O_i(\rr,\khat_i),
\end{equation}
the multi-path form of the master law (\cref{eq:master-law-exact}) with the
activation $g$ defined below. Phase information between paths is discarded.

The coherent levels (7 and 8) act on complex amplitudes. Each path carries a
polarisation-amplitude vector $\psivec_i$, and the body-surface field is the
coherent sum over paths driven by a precoder. The absorbed power density becomes
a Hermitian quadratic form in the precoder $\xx$,
\begin{equation}\label{eq:coherent-form}
  \Sab(\rr) \;=\; \norm{\Gtil(\rr)\,\xx}^2
  \;=\; \xx\herm\,\Qop(\rr)\,\xx,
\end{equation}
with $\Gtil$ the body-surface channel and $\Qop$ the exposure operator, both
developed in \cref{sec:coherent}. \Cref{eq:coherent-form} reduces to
\cref{eq:incoherent-sum} when path phases are randomised.

\subsection{Orthogonal corrections}\label{sec:orthogonal}

Within the incoherent regime the refinements compose because each modifies a
different factor of the per-path term, as shown in
\texttt{composability\_analysis.md}. The general combined per-triangle law is
\begin{equation}\label{eq:combined}
  \Sab \;=\; \underbrace{\Bigl[\Tavg(\mu)+\tfrac{q}{2}\Delta T(\mu)\Bigr]}_{\text{Fresnel}\,+\,\text{polarisation}}
  \;\underbrace{g(\mu,\sigma)}_{\text{activation}}\;S
  \;+\;\underbrace{T_0\,\frac{H}{k}\,g(\mu,\sigma)^2\,S}_{\text{curvature}},
\end{equation}
where the axes are:
\begin{itemize}[leftmargin=1.4em,itemsep=1pt]
  \item \textbf{Fresnel} ($T_0\to\Tavg(\mu)$): the angle-dependent transmission
    of \cref{sec:fresnel}. Modifies the material factor.
  \item \textbf{Polarisation} (the $\tfrac{q}{2}\Delta T$ term): restores the
    TM excess of \cref{sec:pol}, with $q$ the local TM excess. Modifies the
    material factor only.
  \item \textbf{Curvature} (the $H/k$ term): the first-order physical-optics
    correction for a surface of twice-mean-curvature $H$. Adds a perturbative
    $g^2$ term.
  \item \textbf{Diffraction} ($g\in\{\ReLU,\,\mathrm{GeLU},\,\text{Fock}\}$):
    replaces the hard back-face gate $\ReLU(\mu)$ with a physically derived
    smooth transition near the shadow boundary. Modifies the activation only.
  \item \textbf{Inter-body} (the visibility factor $O$): reflections and
    occlusion from other bodies in the scene. Modifies the visibility factor.
\end{itemize}
The corrections are structurally orthogonal: no two modify the same factor, so
any subset can be enabled and the cross-terms between corrections are smaller
than the corrections themselves (for example the polarisation-curvature cross
term is $\sim0.03\%$ for an arm at \SI{28}{\giga\hertz}, three orders of
magnitude below the leading curvature term).

\subsection{The new orthogonal axes}\label{sec:new-axes}

Two axes generalise the original six-level ladder and are the subject of the
later sections. The \textbf{diffraction model} is promoted from a binary
on/off switch to a choice $g\in\{\text{none},\,\text{gelu},\,\text{fock}\}$:
\texttt{none} is the hard $\ReLU$, \texttt{gelu} is the curvature-scaled error
function of \cref{eq:combined}, and \texttt{fock} is the rigorous Fock
creeping-wave model (\cref{sec:fock-local,sec:cmetric}). The
\textbf{inter-body} axis turns the visibility factor $O$ into a coupling that
carries reflected and occluded power between bodies (\cref{sec:interbody}).
Both axes are orthogonal to the Fresnel, polarisation, and curvature axes and
compose with them.

\subsection{Level map}\label{sec:level-map}

\Cref{tab:levels} maps each level to the corrections it enables. Levels 0 and 1
collapse the spatial map to a single aggregate number via the projected area and
absorption directivity. (The viewer's Level 0/1 implementation carries a
documented factor-of-4 caveat in the area normalisation $\Aab$, traced in
\texttt{composability\_analysis.md}: the Cauchy relation
$\langle\Aperp\rangle=A_{\mathrm{total}}/4$ is applied once too often, leaving
the aggregate power four times too small unless $\Aab=A_{\mathrm{total}}$.)

\begin{table}[!htbp]
\centering
\caption{The nine fidelity levels as an orthogonal composition. Fresnel:
  angle-dependent $\Tavg$. Pol: TM excess $\tfrac{q}{2}\Delta T$. Curv:
  physical-optics $H/k$ term. Diff: activation $g$. Coh: coherent quadratic
  form.}
\label{tab:levels}
\begin{tabular}{clccccl}
\toprule
Level & Name & Fresnel & Pol & Curv & Diff & Output \\
\midrule
0 & Worst-case bound      & $T_0$    &   &   & ReLU & scalar bound \\
1 & Aggregate directivity & $T_0$    &   &   & ReLU & scalar power \\
2 & Geometric map         & $T_0$    &   &   & ReLU & $\Sab(\rr)$ \\
3 & Fresnel map           & $\Tavg$  &   &   & ReLU & $\Sab(\rr)$ \\
4 & Polarisation          & $\Tavg$  & \checkmark &   & ReLU & $\Sab(\rr)$ \\
5 & Curvature             & $\Tavg$  & \checkmark & \checkmark & ReLU & $\Sab(\rr)$ \\
6 & Diffraction           & $\Tavg$  & \checkmark & \checkmark & GeLU/Fock & $\Sab(\rr)$ \\
7 & Coherent              & exact $t_s,t_p$ & intrinsic & opt. & opt. & $\norm{\Gtil\xx}^2$ \\
8 & ECBF                  & exact $t_s,t_p$ & intrinsic & opt. & opt. & $\xx\herm\Qop\xx$ \\
\bottomrule
\end{tabular}
\end{table}

In the coherent levels (7 and 8) polarisation is intrinsic: the channel applies
the complex amplitude coefficients $t_s,t_p$ per TE/TM component, so no separate
polarisation switch is needed, and curvature and diffraction enter as real
per-(triangle, path) weights that preserve the quadratic-in-$\xx$ structure.
The remaining sections derive each axis in turn.

% sec_05_polarisation
\section{Polarisation-aware absorption}\label{sec:pol}

\emph{[section pending]}

% sec_06_geometry
\section{Geometry: projection, occlusion, curvature}\label{sec:geom}

With the transmission coefficient fixed at its normal-incidence value $T_0$
(\cref{sec:fresnel}), the absorbed power density of \cref{sec:surface-law}
depends only on the body shape and the incidence direction. The angular factor
\begin{equation}\label{eq:geom-f}
  f(\rr) = \ReLU\!\bigl[\nhat(\rr)\cdot(-\khat)\bigr]
\end{equation}
is purely geometric. It carries the surface-normal field $\nhat(\rr)$ and the
wave direction $\khat$ and nothing else, so millimetre-wave dosimetry reduces to
a problem in integral geometry. This section sets out three pieces of that
geometry: the visibility factor that makes the law non-local on a non-convex
body, the projected-area integral that fixes total absorbed power, and the
principal curvatures that the diffractive gate of \cref{sec:fock-local} needs.

\subsection{Visibility and occlusion}\label{sec:geom-occlusion}

A human body is not convex. Concavities (the armpit, the gap between the legs,
the region behind an ear) let one part of the body shadow another, so the local
value of $\Sab$ is not determined by $\nhat(\rr)$ alone.

\begin{definition}[Occlusion factor]\label{def:occlusion}
For a surface point $\rr$ and an incidence direction $\khat$, the occlusion
factor is
\begin{equation}\label{eq:O-def}
  O(\rr,\khat) = \begin{cases}
    1 & \text{if $\rr$ is visible along $-\khat$ (no other body part blocks it),}\\
    0 & \text{otherwise.}
  \end{cases}
\end{equation}
\end{definition}

The local law becomes
\begin{equation}\label{eq:Sab-occluded}
  \Sab(\rr) = \Sinc\,T_0\,\ReLU\!\bigl[\nhat(\rr)\cdot(-\khat)\bigr]\,O(\rr,\khat).
\end{equation}
Computing $O$ is one ray cast per surface element against the body mesh, the same
operation a renderer performs for hard shadows. For a convex body $O\equiv 1$ and
\cref{eq:Sab-occluded} collapses back to the bare surface law.

The binary factor in \cref{def:occlusion} is the geometric-optics content of
shadowing. It is exact in the limit $\lambda\to 0$ and discontinuous across the
shadow edge. The two smooth, wave-correct refinements of $O$ are the subject of
the next two sections. A body shadowing itself across a local convex edge (its
own terminator) is the Fock gate of \cref{sec:fock-local}. One body part
shadowing another across a finite gap (a hand over the cheek) is the distal
angular-clearance gate of \cref{sec:cmetric}. Both reduce to \cref{eq:O-def} as
$\lambda\to 0$.

\begin{definition}[Exposure fraction]\label{def:eta}
The exposure fraction at $\rr$ is the cosine-weighted fraction of the hemisphere
from which $\rr$ is visible,
\begin{equation}\label{eq:eta-def}
  \eta(\rr) = \frac{1}{\pi}\int_{S^2}
  \ReLU\!\bigl[\nhat(\rr)\cdot(-\khat)\bigr]\,O(\rr,\khat)\,\diff\Omega(\khat).
\end{equation}
\end{definition}

For convex bodies $\eta\equiv 1$. For non-convex bodies $\eta<1$ in concavities
and $\eta=1$ on exposed regions. The exposure fraction is mathematically
identical to the ambient-occlusion factor of computer
graphics~\cite{Zhukov1998}, which imports decades of GPU ray-casting and
bounding-volume-hierarchy machinery into dosimetry.

\subsection{Projected area and the Cauchy bound}\label{sec:geom-cauchy}

Integrating $\Sab$ over the lit surface $\Sigma_+$ (where
$\mu\equiv\nhat\cdot(-\khat)>0$) gives total absorbed power as a single geometric
integral,
\begin{equation}\label{eq:Pabs-Aperp}
  P_{\mathrm{abs}}(\khat) = \Sinc\,T_0\int_{\Sigma_+}\!\mu\,\diff A
  = \Sinc\,T_0\,\Aperp(\khat),
  \qquad
  \Aperp(\khat) \equiv \int_{\Sigma_+}\mu\,\diff A .
\end{equation}
$\Aperp(\khat)$ is the projected (shadow) area of the body normal to $\khat$. For
a convex body it equals the area of the orthogonal projection. For a non-convex
body \cref{eq:Pabs-Aperp} overcounts unless self-occluded elements are removed,
that is unless $\mu$ is replaced by $\mu\,O$; with that correction the identity
holds for any body by a flux argument.

Averaging the projected area over all directions yields the classical Cauchy
result~\cite{Cauchy1841} and its non-convex generalisation through the exposure
fraction.

\begin{theorem}[Generalised Cauchy formula]\label{thm:gen-cauchy}
For any body with surface $\Sigma$ and exposure fraction $\eta$, the
direction-averaged absorbed power is
\begin{equation}\label{eq:gen-cauchy}
  \langle P_{\mathrm{abs}}\rangle
  = \frac{1}{4\pi}\int_{S^2} P_{\mathrm{abs}}(\khat)\,\diff\Omega
  = \Sinc\,T_0\,\frac{\Aab}{4},
  \qquad
  \Aab \equiv \int_\Sigma \eta(\rr)\,\diff A .
\end{equation}
\end{theorem}

\begin{proof}
Exchange the order of integration (Fubini, the integrand is non-negative) and use
the rotational symmetry of the inner integral,
$\int_{S^2}\ReLU[\nhat\cdot(-\khat)]\,\diff\Omega = \pi$, independent of $\nhat$.
With the occlusion factor restricting the surface integral to visible elements,
$\langle P_{\mathrm{abs}}\rangle = (\pi/4\pi)\,\Sinc\,T_0\,\Aab = \Sinc\,T_0\,\Aab/4$.
\end{proof}

The absorption area $\Aab\le A$ collapses all of self-shadowing into one scalar;
for convex bodies $\eta\equiv 1$ and $\Aab=A$, recovering $\langle\Aperp\rangle =
A/4$. Replacing $T_0$ with the exact averaged transmission $\Tavg$
(\cref{sec:fresnel}) makes \cref{eq:gen-cauchy} exact at any frequency.

\subsection{Principal curvatures and the Fock radius}\label{sec:geom-curvature}

The geometric law treats each surface element as locally flat. The wave
corrections do not: the curvature near the shadow edge sets the width of the
diffractive penumbra. The relevant quantity is not the mean curvature but the
normal curvature in the plane of incidence.

At a smooth surface point the shape operator has two real eigenvalues, the
principal curvatures $\kappa_1\ge\kappa_2$, with orthonormal principal directions
$\bm{e}_1,\bm{e}_2$ in the tangent plane (sign convention: $\kappa>0$ for an
outward-bulging convex surface). Their sum is twice the mean curvature,
$H=\kappa_1+\kappa_2$, the quantity that enters the physical-optics curvature
term $\Sab \propto \ReLU(\mu)\,[1+\mu H/k]$ used at the curvature level of the
grid (\cref{sec:grid}). AEGIS estimates $\kappa_1,\kappa_2,\bm{e}_1$ per triangle
from a local quadric (osculating-jet) fit over a centroid neighbourhood,
validated on analytic sphere and cylinder meshes to a few percent.

The normal curvature in an arbitrary tangent direction follows from Euler's
theorem. Let $\varphi$ be the angle between $\bm{e}_1$ and the incidence ray
projected into the tangent plane. Then
\begin{equation}\label{eq:euler}
  \kappa_t(\varphi) = \kappa_1\cos^2\varphi + \kappa_2\sin^2\varphi,
\end{equation}
and the in-incidence-plane radius of curvature is $\Rfock = 1/\kappa_t$. This
$\Rfock$, not $2/H$, is the radius the Fock detour parameter of
\cref{sec:fock-local} requires.

\begin{proposition}[The Fock radius is the in-plane normal curvature]\label{prop:fock-radius}
Let a ray strike a circular cylinder of radius $R$ normal to its axis. The
in-incidence-plane Fock radius is $\Rfock = R$, whereas the mean-curvature
surrogate gives $2/H = 2R$, twice too large. The two disagree by a factor of $2$,
and the diffractive penumbra width $\propto(kR)^{-1/3}$ they predict differs by
$2^{1/3}\approx 1.26$.
\end{proposition}

\begin{proof}
For a cylinder of radius $R$ the principal curvatures are $\kappa_1=1/R$ (around
the cross-section) and $\kappa_2=0$ (along the axis), so $H=\kappa_1+\kappa_2=1/R$
and $2/H=2R$. A ray normal to the axis projects, in the tangent plane, onto the
cross-sectional principal direction $\bm{e}_1$, so $\varphi=0$ and
\cref{eq:euler} gives $\kappa_t=\kappa_1=1/R$, hence $\Rfock=R$. The penumbra
width scales as $(k\Rfock)^{-1/3}$ (\cref{sec:fock-local}); substituting $2R$ for
$R$ rescales it by $2^{1/3}$.
\end{proof}

The cylinder is a primary diffraction oracle (\cref{sec:validation}), so the
factor-of-two error is not academic: it would corrupt the very case the gate is
validated against. For a sphere $\kappa_1=\kappa_2=1/R$ and \cref{eq:euler} gives
$\Rfock=R$ for every direction, as it must by isotropy. On a flat face
$\kappa_t\to 0$ and $\Rfock\to\infty$, which returns the gate to a sharp ReLU.

% sec_07_fock_local
\section{Local diffraction: the Fock gate}\label{sec:fock-local}

\emph{[section pending]}

% sec_08_distal_cmetric
\section{Distal occlusion: the angular-clearance metric}\label{sec:cmetric}

\emph{[section pending]}

% sec_09_interbody
\section{Inter-body reflection}\label{sec:interbody}

The master surface law of \cref{sec:surface-law} captures first-bounce
absorption only. At each illuminated point a fraction $1-T_0$ of the
incident power is reflected rather than absorbed, and on a non-convex
body part of that reflected power can re-illuminate another body part:
across the gap between the legs, between an arm and the torso, or under
the chin. This section formalises the single specular recapture as an
optional second-order correction. It is off by default. The production
law omits it because the validated body-averaged enhancement stays
inside the diffraction and tissue-property error bars.

\subsection{Reflected power and the specular direction}

At a front-facing point $\rr$ with incidence cosine
$\mu=\nhat\cdot(-\khat)$ the Fresnel reflection coefficients are the
siblings of the transmission coefficients of \cref{sec:fresnel},
\begin{equation}\label{eq:r-coeffs}
  r_s(\mu)=\frac{\mu-\zeta}{\mu+\zeta},\qquad
  r_p(\mu)=\frac{\ntilde^2\mu-\zeta}{\ntilde^2\mu+\zeta},\qquad
  \zeta=\sqrt{\ntilde^2-1+\mu^2},\ \ \RE\zeta>0,
\end{equation}
and they satisfy the polarisation-wise energy identity
$T_s+\abs{r_s}^2=1$, $T_p+\abs{r_p}^2=1$ at the lossy interface.

\begin{definition}[Reflected power fraction]\label{def:rho-refl}
With TE and TM power weights $w_s,w_p$ ($w_s+w_p=1$) carried from the
incident polarisation as in \cref{sec:pol}, the locally reflected
fraction of the incident power density is
\begin{equation}\label{eq:rho-refl}
  \rho_R(\rr)=\abs{r_s(\mu)}^2 w_s+\abs{r_p(\mu)}^2 w_p
            =1-T_0(\rr),
\end{equation}
the Fresnel reflectance. The specularly reflected ray leaves along
\begin{equation}\label{eq:k-refl}
  \khat_{\mathrm{refl}}=\khat-2(\khat\cdot\nhat)\,\nhat
                       =\khat+2\mu\,\nhat,
\end{equation}
the mirror of $\khat$ about the tangent plane.
\end{definition}

\subsection{One-bounce recapture}

The single specular recapture casts the reflected ray
\cref{eq:k-refl} from every lit triangle and reuses the visibility
acceleration structure of \cref{sec:geom} to test whether it
strikes another body triangle before escaping. When it does, the
reflected power density $\rho_R\,\Sinc$ is deposited at the hit point,
weighted by its own incidence cosine, and added to the direct
first-bounce term there. This is the $\mathtt{inter\_body=specular1}$
axis: a single optional pass, orthogonal to the diffraction and
polarisation flags, that never feeds back into the visibility bake.

To bound the cumulative effect of all bounces, let $f(\rr)$ denote the
fraction of power reflected from $\rr$ that strikes the body again.
Summing the geometric series of successive bounces enhances the
absorbed power at a point by
\begin{equation}\label{eq:radiosity-C}
  C(f)=\frac{1}{1-\Tbar_R\,f},\qquad \Tbar_R=1-\Tbar\approx 0.46,
\end{equation}
with $\Tbar_R$ the flux-weighted reflectance. Under diffuse
(Lambertian) reflection the recaptured fraction is bounded by the
cosine-weighted hemisphere blocked by other body parts,
\begin{equation}\label{eq:f-bound}
  f(\rr)\le 1-\Omega(\rr),
\end{equation}
where $\Omega$ is exactly the ambient-occlusion exposure fraction of
\cref{sec:geom}. For a convex body every reflected ray escapes,
$f\equiv 0$, and the first-bounce law is exact.

\begin{proposition}[Self-compensation]\label{prop:self-comp}
The enhancement $C$ is largest precisely where $\Omega$ is smallest. A
deep concavity sees little direct power but recaptures a large share of
the reflected power, so the product $\Omega\,C$ stays close to $\Omega$.
By \cref{eq:f-bound} a point with $\Omega=0.3$ admits $f\le 0.7$, hence
$C\le 1/(1-0.46\times 0.7)\approx 1.47$ and $\Omega\,C\le 0.44$,
recovering less than a fifth of the $70\%$ shadow deficit.
\end{proposition}

\subsection{Energy bookkeeping}

The recapture is an exchange of power between body parts inside the
incidence-operator balance derived for the coherent regime in
\cref{sec:coherent}. Writing $\Qop$ for the absorption operator and
$\Qop_{\mathrm{ref}}$ for its reflection complement (the Gram of the
specularly scattered surface field), the landed power partitions as
\begin{equation}\label{eq:op-balance}
  \Qop_{\mathrm{inn}}=\Qop+\Qop_{\mathrm{ref}},
\end{equation}
where $\Qop_{\mathrm{inn}}$ maps a precoder to the power that lands on
the body irrespective of whether it is then absorbed or scattered
away. The single specular recapture promotes a slice of
$\Qop_{\mathrm{ref}}$ back into $\Qop$ to second order in the body's
own albedo $\bar\alpha=\operatorname{tr}\Qop_{\mathrm{ref}}/
\operatorname{tr}\Qop_{\mathrm{inn}}$, leaving evanescent surface
waves, contour creeping, and triple-and-higher bounces in the
residual.

\subsection{Validated bounds}

Measured on the Thelonious phantom (\num{23826} triangles) over a sweep
of incidence directions, the correction is small but the monograph's
original specular estimate was optimistic.

\begin{remark}[What the data say]\label{rem:interbody-data}
Three findings hold across every concavity region.
(i) The body-averaged enhancement is $C\approx 1.05$ to $1.08$, at most
$8\%$, well below the diffraction (${\sim}10\%$) and tissue-property
(${\sim}10\%$) uncertainties, which justifies omitting reflection from
the default law.
(ii) The specular recapture is \emph{comparable} to the diffuse bound,
not one third of it: in a concavity the mirror direction points back
at the opposing surface, giving
$f_{\mathrm{spec}}\approx 1.08\,f_{\mathrm{diffuse}}$, so the specular
enhancement matches the diffuse $C$ rather than the near-unity value
once assumed.
(iii) The deepest concavity reaches a local $C\approx 1.85$. This
matters for spatial maps in tight gaps, not for compliance, which is
set by the most-illuminated surface where $f\to 0$.
\end{remark}

The correction is therefore retained as the optional
$\mathtt{specular1}$ pass for spatial-map fidelity and recorded in the
error budget of \cref{sec:validation}, but it is excluded from the
default equations.

% sec_10_nearfar
\section{Near and far field}\label{sec:nearfar}

The surface law of \cref{sec:surface-law} is stated for a plane wave: a
single incidence direction $\khat$ and a single incident power density
$\Sinc$ shared by every surface point. That is the far-field limit of a
source receding to infinity. This section places it in a distance
hierarchy and gives the point-source generalisation that the
differentiable design loop and the coherent phase of \cref{sec:coherent}
both build on.

\subsection{Three distance regimes}

Let $d(\rr)=\abs{\rr-\rr_{\mathrm{src}}}$ be the distance from the
source to a surface point and $L$ the largest body dimension.

\begin{definition}[Reactive near field, $d<\lambda/2\pi$]\label{def:reactive}
The field is dominated by non-propagating evanescent components and the
plane-wave surface law does not apply. At \SI{28}{\giga\hertz} this
boundary is \SI{1.7}{\milli\metre} from the source, so only sources in
direct skin contact fall in this regime.
\end{definition}

\begin{definition}[Locally plane-wave regime, $d>3\lambda$]\label{def:lpw}
The reactive field has decayed and the wavefront curvature is small
enough over a single triangle that the illumination is locally plane
but spatially non-uniform. The surface law holds pointwise with inputs
that vary across the body,
\begin{equation}\label{eq:Sab-point}
  \Sab(\rr)=\frac{P_t\,G\!\bigl(\khat(\rr)\bigr)}{4\pi\,d(\rr)^2}\,
            T_0\,\ReLU\!\bigl[\nhat(\rr)\cdot(-\khat(\rr))\bigr],
\end{equation}
where $P_t$ is the radiated power, $G$ the antenna gain pattern, and
the per-point incidence direction is
\begin{equation}\label{eq:khat-r}
  \khat(\rr)=\frac{\rr-\rr_{\mathrm{src}}}{d(\rr)}
            =\widehat{\rr-\rr_{\mathrm{src}}}.
\end{equation}
The pseudo-Brewster compensation and the ReLU cosine gate carry over
unchanged. The cost is $O(M)$ per source, the same as a plane wave.
\end{definition}

\begin{definition}[Far field, $d>5L$]\label{def:farfield}
Wavefront curvature and distance variation are negligible over the
whole body: $d(\rr)\to d_0$ and $\khat(\rr)\to\khat_0$ for every
surface point, recovering the plane-wave law with
$\Sinc=P_t\,G(\khat_0)/(4\pi d_0^2)$. The approximation error scales as
$(L/2d)^2$: below $1\%$ at $d>\SI{4}{\metre}$ and below $3\%$ at
$d\approx\SI{2}{\metre}$ for a standing human illuminated from the
front.
\end{definition}

\subsection{Point source and the view factor}

Integrating \cref{eq:Sab-point} over the visible front-facing surface,
the factor $\mu(\rr)\,\diff A/d(\rr)^2$ is the differential solid angle
$\diff\Omega_s$ subtended by a surface element at the source. For an
isotropic source this collapses to a view-factor law,
\begin{equation}\label{eq:Pabs-solidangle}
  P_{\mathrm{abs}}=\frac{P_t\,T_0}{4\pi}\,\Omega_{\mathrm{body}}(\rr_{\mathrm{src}}),
  \qquad
  \Omega_{\mathrm{body}}=\int_{\Sigma_+}\frac{\mu(\rr)}{d(\rr)^2}\,\diff A,
\end{equation}
so absorbed power equals $T_0$ times the radiated power times the view
factor of the body seen from the source. As $\rr_{\mathrm{src}}$ recedes,
$\Omega_{\mathrm{body}}\to\Aperp(\khat_0)/d_0^2$ and
$P_{\mathrm{abs}}\to\Sinc\,T_0\,\Aperp$; as the source approaches the
skin, $\Omega_{\mathrm{body}}\to 2\pi$ and
$P_{\mathrm{abs}}\to P_t\,T_0/2$. A directive antenna carries its gain
pattern into the integrand as a weight.

\subsection{The wavefront factor}

The finite distance $d_1$ to the last scatter enters the rest of the
framework through one geometric factor, in two places.

First, it sets the diffraction penumbra width. The Fock detour scale of
\cref{sec:fock-local} acquires the source-distance correction
$w_{\mathrm{nf}}$ of \cref{sec:cmetric}: a converging or diverging
wavefront widens or narrows the geometric shadow transition relative to
the plane-wave value, and the gate is evaluated at the corrected detour.

Second, it sets the coherent phase. In the coherent regime of
\cref{sec:coherent} the plane-wave phase $e^{-ik_0\khat_n\cdot\rr}$ is
replaced by the near-field phase kernel
\begin{equation}\label{eq:NF-kernel}
  \varphi_n(\rr)=\frac{R_n^{(0)}}{R_n(\rr)}\,
                 e^{-ik_0\bigl(R_n(\rr)-R_n^{(0)}\bigr)},
  \qquad R_n(\rr)=\abs{\rr-\rr_n^{\mathrm{scat}}},
\end{equation}
referenced to a fixed body point $\rr^{(0)}$ so that
$\varphi_n(\rr^{(0)})=1$. Unlike the plane-wave phase diagonal, this
kernel is diagonal but \emph{not} unitary: it carries both the
$e^{-ik_0 R_n}$ phase and the $1/R_n$ amplitude decay. With this
substitution the channel factorisation, the exposure operator, and the
exposure-constrained precoder of \cref{sec:coherent} all proceed
verbatim, because the derivation uses only the diagonal structure of
the kernel, not its specific form. The coherent hotspot shape changes:
the constant-wave-vector Fourier cells of the far field deform into
confocal lobes whose foci are the last-scatter points of the
interfering path pair.

\subsection{Differentiability in the source position}

Both $d(\rr)$ and $\khat(\rr)$ are smooth functions of the source
position $\rr_{\mathrm{src}}$, and the ReLU gate is differentiable
almost everywhere with a defined subgradient at the terminator. The
point-source law \cref{eq:Sab-point} and the view-factor integral
\cref{eq:Pabs-solidangle} are therefore differentiable in
$\rr_{\mathrm{src}}$, which is what lets the design loop optimise source
placement by gradient descent rather than by sweep. The view factor
also admits an antenna-agnostic factorisation: expanding the gain in
spherical harmonics, $G(\hat u)=\sum_{lm}g_{lm}Y_{lm}(\hat u)$, and
exchanging the sum with the surface integral gives
\begin{equation}\label{eq:Gamma-lm}
  P_{\mathrm{abs}}=\frac{T_0}{4\pi}\sum_{lm}g_{lm}\,\Gamma_{lm}(\rr_{\mathrm{src}}),
  \qquad
  \Gamma_{lm}=\int_{\Sigma_+}\frac{\mu(\rr)\,Y_{lm}\!\bigl(\khat(\rr)\bigr)}
                                  {d(\rr)^2}\,O(\rr,\khat)\,\diff A,
\end{equation}
so the body response $\Gamma_{lm}$ depends on the source position but
not the antenna pattern. Any antenna at a fixed position is then
evaluated by one dot product, the acceleration that makes
phone-scale dosimetry over many antenna orientations tractable.

% sec_11_coherent
\section{Coherent dosimetry: the exposure operator}\label{sec:coherent}

\emph{[section pending]}

% sec_12_pssar
\section{The peak-spatial-SAR bridge}\label{sec:pssar}

% Local symbols not defined in the shared preamble. \providecommand is safe:
% it defines only if the master has not already introduced them.
\providecommand{\SAR}{\mathrm{SAR}}
\providecommand{\psSAR}{\mathrm{psSAR}_{10\mathrm{g}}}
\providecommand{\rhom}{\rho_{\mathrm{m}}}
\providecommand{\Tzero}{T_{0}}

Everything so far produces $\Sab$, the absorbed power density on the body
surface. Above 6\,GHz that is the quantity ICNIRP restricts directly. Below
6\,GHz the regulated metric is instead the peak spatial-average specific
absorption rate over a 10\,g cube of tissue, $\psSAR$, a volumetric quantity.
This section bridges the two. Under a homogeneous lossy half-space at each
surface point, the exponential depth decay of the internal fields collapses the
three-dimensional cube average to a two-dimensional surface operation on
$\Sab$. The result is a closed-form $\psSAR$ from the same surface map, a
one-line dominance bound that certifies $\psSAR$ compliance from any $\Sab$
calculation, and a correction to the standardized averaging cube that the naive
projection argument gets wrong.

\subsection{Internal fields from surface absorbed power density}
\label{sec:pssar-internal}

At a surface point $\rr$ the transmitted field in tissue decays into depth $z$
(measured along the inward normal) as $e^{-\alpha z}$ in amplitude, where
$\alpha = k_0\kappa$ is the normal-incidence field decay rate and
$\ntilde = n - \mathrm{i}\kappa$. Because $\abs{\ntilde}\gg 1$ for biological
tissue (for example $\abs{\ntilde} = 4.84$ for skin at 28\,GHz), Snell's law
compresses all refraction angles to within roughly $10^\circ$ of normal, so
$\alpha$ varies by at most 2.2\% across $\theta_i\in[0^\circ,85^\circ]$. We
adopt the normal-incidence value throughout (the universal depth-coupling
approximation).

Integrating the ohmic loss density $\tfrac{\sigma}{2}\abs{\EE}^2$ over depth and
matching the total to $\Sab$ fixes the surface field intensity, hence the full
depth profile.

\begin{lemma}[Depth-resolved field and SAR]\label{lem:depth-profile}
Under the homogeneous half-space and universal-$\alpha$ approximations,
\begin{equation}\label{eq:E-depth}
  \abs{\EE(\rr,z)}^2 = \frac{4\alpha}{\sigma}\,\Sab(\rr)\,e^{-2\alpha z},
  \qquad
  \SAR(\rr,z) = \frac{2\alpha}{\rhom}\,\Sab(\rr)\,e^{-2\alpha z},
\end{equation}
with $\rhom$ the tissue mass density. The surface value is
$\SAR(\rr,0) = (2\alpha/\rhom)\,\Sab(\rr)$, and
$\rhom\int_0^\infty \SAR(\rr,z)\,\diff z = \Sab(\rr)$ by construction.
\end{lemma}

The relation $\abs{\EE(\rr,0)}^2 = (4\alpha/\sigma)\Sab$ is exact for TE
polarisation. For TM there is a multiplicative correction
$C_p = 1 + \sin^2\theta/\abs{\xi}^2$ with
$\xi = \sqrt{\ntilde^2 - 1 + \mu^2}$, bounded by $4.5\%$ for skin at 28\,GHz and
below $2\%$ for $\theta\le 60^\circ$. The characteristic SAR depth is
$1/(2\alpha) \approx 0.48$\,mm for skin at 28\,GHz, so the absorption is
concentrated in a layer far thinner than the cube.

\subsection{The 10\,g average as a surface functional}
\label{sec:pssar-functional}

\begin{definition}[Peak spatial-average SAR]\label{def:pssar}
\begin{equation}\label{eq:pssar-def}
  \psSAR = \sup_{\Omega}\;\frac{1}{m}\int_\Omega \frac{\sigma}{2}\,\abs{\EE}^2\,\diff V,
  \qquad m = \rhom\operatorname{Vol}(\Omega) = 10\;\mathrm{g}.
\end{equation}
ICNIRP 1998 admitted any contiguous tissue region $\Omega$; ICNIRP 2020
restricts $\Omega$ to a cube, and IEC/IEEE 62704-1 fixes the side at
$L = (m/\rhom)^{1/3} = 21.5$\,mm.
\end{definition}

At mmWave the skin depth (0.3 to 0.5\,mm) is far smaller than the minimum body
radius of curvature, so the optimal averaging volume is a thin
surface-conforming slab: a connected patch $A\subset\Sigma$ extruded to depth
$d = m/(\rhom A)$. Substituting \cref{eq:E-depth} and performing the depth
integral analytically yields the depth-efficiency function
\begin{equation}\label{eq:h-def}
  h(u) = \frac{1 - e^{-u}}{u}, \qquad u = 2\alpha d = \frac{2\alpha m}{\rhom A},
\end{equation}
which is strictly decreasing, concave, and satisfies $h(u)\in(0,1]$ with
$h(0^+)=1$. It is the fraction of surface SAR that survives mass-averaging to
depth $d$.

\begin{theorem}[psSAR from surface APD, relaxed geometry]\label{thm:pssar-slab}
Under the homogeneous half-space and universal-$\alpha$ approximations,
\begin{equation}\label{eq:pssar-slab}
  \psSAR = \frac{2\alpha}{\rhom}\;
    \sup_{\substack{A\subset\Sigma \\ A\;\mathrm{connected}}}
    \left[\langle\Sab\rangle_A\;h\!\Bigl(\frac{2\alpha m}{\rhom A}\Bigr)\right],
\end{equation}
where $\langle\Sab\rangle_A = A^{-1}\int_A \Sab\,\diff A$ is the area average over
the patch.
\end{theorem}

The supremum balances two effects. A small patch captures the hotspot (large
$\langle\Sab\rangle_A$) but extends deep into tissue where SAR has decayed
(small $h$); a large patch keeps absorption near the surface ($h\approx 1$) but
dilutes $\langle\Sab\rangle_A$ with low-valued regions. The maximiser balances
hotspot capture against depth dilution. \Cref{eq:pssar-slab} relaxes the cube
shape, so it upper-bounds the standardized cube result of
\cref{sec:pssar-cube}.

For a sphere of radius $R$ illuminated along a principal axis, the patch is a
cap of half-angle $\theta_0$ and \cref{eq:pssar-slab} reduces to a single
transcendental equation in the dimensionless ratio
$\lambda = \alpha m/(\rhom\pi R^2)$. At $R = 15$\,cm, 28\,GHz,
$\Sinc = 10\;\mathrm{W/m^2}$ the optimum is a thin $0.21$\,mm slab spanning
$\approx 479\;\mathrm{cm}^2$ with $\psSAR \approx 7.59\;\mathrm{W/kg}$,
reflecting the exponential concentration of SAR near the surface.

\subsection{The APD-dominance bound}
\label{sec:pssar-dominance}

The two ICNIRP local basic restrictions, $\psSAR$ below 6\,GHz and $\Sab$ above
6\,GHz, are tied together by one inequality from energy conservation. The cube
can hold no more absorbed power than enters through its body-surface footprint.

\begin{proposition}[APD dominates psSAR]\label{prop:apd-dom}
For an axis-aligned 10\,g cube placed per IEC/IEEE 62704-1 on a planar
three-layer body, with the cube axis aligned to the local outward normal
$\nhat$,
\begin{equation}\label{eq:apd-dom}
  \psSAR \le \frac{\Sab}{\rhom L},
\end{equation}
where $\Sab$ is the absorbed power density at the cube footprint. Equality holds
in the strong-attenuation limit $\alpha L\to\infty$.
\end{proposition}

\begin{proof}
$\Sab$ is the absorbed flux entering tissue per unit footprint area, so the
absorbed power inside the cube satisfies
$P_{\mathrm{abs}} \le \Sab\,A$ with footprint $A = L^2$ for the aligned cube.
Cube mass is $m = \rhom L^3$. Dividing,
$\psSAR = P_{\mathrm{abs}}/m \le \Sab L^2/(\rhom L^3) = \Sab/(\rhom L)$. The
share of entering flux deposited inside the cube is the capture fraction
$C(\alpha L) = 1 - e^{-2\alpha L}$, which $\to 1$ as $\alpha L\to\infty$.
\end{proof}

With $\rhom = 1000\;\mathrm{kg/m^3}$ and $L = 21.5$\,mm,
$\rhom L = 21.5\;\mathrm{kg/m^2}$, so the ICNIRP $\Sab$ basic restriction of
$10\;\mathrm{W/m^2}$ implies $\psSAR \le 0.465\;\mathrm{W/kg}$. This is below
the head-and-trunk $\psSAR$ restriction ($2\;\mathrm{W/kg}$) by a factor of
$4.3$ and below the limb restriction ($4\;\mathrm{W/kg}$) by a factor of $8.6$.
A cube whose axes do not align with $\nhat$ intercepts a tilted body plane; the
worst-case face-diagonal cut carries a $\sqrt{2}$ prefactor,
\begin{equation}\label{eq:apd-dom-gen}
  \psSAR \le \frac{\sqrt{2}\,\Sab}{\rhom L},
\end{equation}
giving $\psSAR \le 0.658\;\mathrm{W/kg}$ at the $\Sab$ limit, still $3.0\times$
below the head-and-trunk restriction. The bound is tight to within the capture
fraction, which exceeds $0.98$ above 5\,GHz on skin and falls to $0.67$ at
0.9\,GHz where the penetration depth exceeds the cube side
(\cref{tab:pssar-capture}). A Monte Carlo 62704-style search with a three-layer
skin/fat/muscle model on the Thelonious phantom across 1 to 28\,GHz returns a
median bound-to-numerical ratio of $1.30$ and a maximum of $1.46$, confirming
\cref{eq:apd-dom} in realistic anatomy.

\begin{table}[ht]
\centering
\caption{Field decay rate $\alpha$ on skin and capture fraction
$C(\alpha L) = 1 - e^{-2\alpha L}$ for the $L = 21.5$\,mm cube. Above 5\,GHz the
dominance bound is exact to a few percent.}
\label{tab:pssar-capture}
\begin{tabular}{@{}rccc@{}}
\toprule
$f$ [GHz] & $1/(2\alpha)$ [mm] & $2\alpha L$ & $C$ \\
\midrule
0.9  & 19.9 & 1.1  & 0.67 \\
2.45 & 11.3 & 1.9  & 0.85 \\
5    &  5.3 & 4.1  & 0.98 \\
10   &  1.9 & 11.3 & 1.00 \\
28   & 0.46 & 46.7 & 1.00 \\
60   & 0.24 & 89.9 & 1.00 \\
\bottomrule
\end{tabular}
\end{table}

The operational consequence is that the same surface integral that delivers
$\Sab$ pointwise and whole-body absorbed power through the generalised Cauchy
identity also certifies $\psSAR$ compliance, with no volumetric cube search. The
6\,GHz switch between basic restrictions is a regulatory convention, not a
computational necessity. The bound does not say $\psSAR$ should be retired:
$\psSAR$ encodes deeper-tissue heating through the cube depth that $\Sab$ does
not, and its limit values were calibrated against local temperature-rise models
that include this contribution.

\subsection{The standardized cube and the Cauchy-projection critique}
\label{sec:pssar-cube}

A tempting shortcut treats the cube as an opaque object suspended in the
wavefront, so that absorbed power equals $\Sinc\Tzero$ times the cube's
projected area $A_\perp(\khat) = L^2(\abs{k_x}+\abs{k_y}+\abs{k_z})$, giving the
Cauchy-projection corollary
\begin{equation}\label{eq:cauchy-form}
  \langle\SAR\rangle_{\mathrm{cube}}
  \overset{?}{=} \frac{\Sinc\Tzero}{\rhom L}\,
    \bigl(\abs{k_x}+\abs{k_y}+\abs{k_z}\bigr).
\end{equation}
A face-by-face flux audit of an actual 62704-valid cube shows
\cref{eq:cauchy-form} is wrong outside narrow regimes. For a cube placed
face-flush on a smooth body, five of the six faces are immersed in tissue and
intercept no fresh flux; the tissue strips on the side faces channel a Poynting
inflow of at most $\sim 10^{-4}\,\Sinc\Tzero L^2$ per face at 28\,GHz, because a
photon entering high-index tissue ($\abs{\ntilde}\gtrsim 4$) refracts to near
normal and is absorbed within $1/(2\alpha)\approx 0.5$\,mm, far short of the
$21.5$\,mm cube. All absorbed power therefore enters through the single
air-facing face, and a direct surface accounting gives
\begin{equation}\label{eq:cube-correct}
  P_{\mathrm{abs}}^{\mathrm{cube}}
  = \int_{\Sigma\cap\mathcal{C}}\Sab\,\diff A
  = \Sinc\Tzero L^2\cos\theta_i,
  \qquad
  \langle\SAR\rangle_{\mathrm{cube}} = \frac{\Sab(\rr)}{\rhom L},
\end{equation}
proportional to $\cos\theta_i$, not to $\abs{k_x}+\abs{k_y}+\abs{k_z}$.

\begin{remark}[Where the projection breaks]\label{rem:cauchy-break}
The Cauchy projection counts the full geometric shadow of the cube as if the
body filled it on every forward-facing side. A 62704-valid cube is mostly inside
the body, so the control surface is the body patch $\Sigma\cap\mathcal{C}$, not
the cube boundary $\partial\mathcal{C}$. The projection is exact only when the
body genuinely fills the shadow on multiple faces: a sharp convex corner,
axis-aligned incidence (where $\abs{k_x}+\abs{k_y}+\abs{k_z}=1=\cos\theta_i$),
or a slender feature engulfed by the cube. None of these is the dominant
plane-wave-on-torso compliance scenario.
\end{remark}

A numerical 62704-style optimisation on Thelonious (four frequencies, two
polarisations, five incidence directions, 280 valid records) makes the failure
quantitative. The Cauchy formula overpredicts $\psSAR$ by a median of
$3.1\times$ and diverges to $10$ to $50\times$ at grazing incidence, and it is
non-monotone in $\theta_i$ (peaking near $45^\circ$) contrary to the physics.
The corrected cosine law of \cref{eq:cube-correct} stays within $\sim 50\%$ of
numerical, with the residual driven by the tilt $\gamma$ between body normal and
cube axis (\cref{tab:pssar-cube}): essentially exact for $\gamma\le 15^\circ$,
degrading monotonically to $+50\%$ at $\gamma\approx 50^\circ$ as the
axis-aligned cube samples a depth distribution biased away from the surface.

\begin{table}[ht]
\centering
\caption{Median ratio of the corrected cosine law \cref{eq:cube-correct} to the
numerical 62704 maximum on Thelonious, binned by the tilt $\gamma$ between body
normal and nearest cube axis. The Cauchy projection \cref{eq:cauchy-form}
instead overpredicts by a median $3.1\times$ across the same sweep.}
\label{tab:pssar-cube}
\begin{tabular}{@{}rccccc@{}}
\toprule
$\gamma$ & 2.45\,GHz & 10\,GHz & 28\,GHz & 60\,GHz & median \\
\midrule
5--15$^\circ$  & 1.16 & 0.94 & 0.95 & 0.95 & 0.96 \\
15--25$^\circ$ & 1.26 & 1.03 & 1.19 & 1.24 & 1.20 \\
25--35$^\circ$ & 1.35 & 1.09 & 1.33 & 1.40 & 1.33 \\
35--45$^\circ$ & 1.43 & 1.14 & 1.39 & 1.42 & 1.40 \\
45--55$^\circ$ & 1.46 & 1.16 & 1.45 & 1.50 & 1.46 \\
\bottomrule
\end{tabular}
\end{table}

\subsection{Domain of validity}
\label{sec:pssar-domain}

The bridge rests on three approximations: the universal-$\alpha$ decay rate
(error $\le 2.2\%$ at 28\,GHz), the flat-slab volume approximation (curvature
correction $\mathcal{O}(d/R)\sim 0.5\%$ for torso and limbs), and the Fresnel
half-space model (valid when the skin depth is much smaller than the body
dimensions). All three hold in the mmWave surface-dominated regime, where the
absorption layer $1/(2\alpha)$ collapses to a fraction of a millimetre and the
cube average becomes a surface average. Below a few GHz the capture fraction
degrades, layered standing-wave effects appear, and the half-space picture
weakens; there the full depth integral or a layered transfer-matrix model is
needed. For thin or folded anatomy (ears, fingers) where the body surface
doubles back inside a single cube, the projected-area shortcut and the
single-interface assumption both fail, and the general surface integral
$\int_{\Sigma\cap\mathcal{C}}\Sab\,\diff A$ over all illuminated patches must be
used directly.

% sec_13_validation
\section{Validation against exact references}\label{sec:validation}

\emph{[section pending]}


\end{document}
